\section*{अध्याय \ref{ch:b2:waves-in-media} --- प्लाज़्मा, चालक और परावैद्युत में विद्युत्चुंबकीय तरंगें}

\begin{solution}{exo:b2:waves-in-media:1}
$f_p \approx 9\sqrt n$ Hz: \qty{2.8}{MHz}, \qty{9}{MHz}, \qty{9}{GHz}, \qty{2.6e15}{Hz}
($\lambda_p = \qty{115}{nm}$), \qty{2.8}{kHz}। \qty{1}{MHz} पर AM: केवल
अंतरतारकीय अंतरिक्ष से। FM: आयनमंडल से, आवरण से नहीं। GPS:
आवरण को छोड़ सब कुछ। प्रकाश: ताँबे को छोड़ सब कुछ।
\end{solution}

\begin{solution}{exo:b2:waves-in-media:2}
ताँबा: \qty{9.2}{mm}, \qty{2.1}{mm}, \qty{65}{\micro m}, \qty{2.1}{\micro m}; \qty{1}{MHz} पर
ऐलुमिनियम: \qty{84}{\micro m}; समुद्री जल: \qty{10}{kHz} पर \qty{2.3}{m}, \qty{1}{GHz} पर
\qty{7}{mm}; मिट्टी: \qty{1}{MHz} पर \qty{5}{m}। \qty{10}{m} पर किसी पनडुब्बी को चाहिए $\delta
\gtrsim \qty{10}{m}$: अर्थात् लगभग \qty{500}{Hz} से नीचे।
\end{solution}

\begin{solution}{exo:b2:waves-in-media:3}
$n = 1.528$, $1.515$, $1.509$; $v = 1.963$, $1.980$, $\qty{1.987e8}{m/s}$; $n_g = n +
2B/\lambda^2 = 1.545$। अपवर्तन कोण $\arcsin(0.707/n)$: \qty{27.57}{\degree} और
\qty{27.93}{\degree}: अर्थात् \qty{0.36}{\degree} का फैलाव।
\end{solution}

\begin{solution}{exo:b2:waves-in-media:4}
$f_p = \qty{4.0}{MHz}$: \qty{5}{MHz} पार जाता है; \qty{2}{MHz} परावर्तित होता है, और
$c/\sqrt{\omega_p^2 - \omega^2} = \qty{14}{m}$ तक घुसता है। $n = \qty{1.5e12}{m^{-3}}$ के लिए: \qty{11}{MHz}।
घनत्व सूर्य के आयनकारी फ्लक्स के पीछे चलता है --- घंटा, ऋतु, और
ग्यारह-वर्षीय चक्र --- अतः सबसे ऊँची काम की आवृत्ति उनके साथ खिसकती है।
\end{solution}

\begin{solution}{exo:b2:waves-in-media:5}
(क) $1/\gamma\pi a^2 = \qty{5.3}{m\ohm/m}$। (ख) $\delta = \qty{65}{\micro m}$; $R \approx 1/\gamma2\pi a\delta
= \qty{41}{m\ohm/m}$, अर्थात् दिष्ट का $7.7$ गुना। (ग) $\delta = \qty{6.5}{\micro m}$, \qty{0.41}{\ohm/m},
$77$ गुना। (घ) त्वचा के हावी होते ही $R \propto \sqrt f$ उतनी ही तेज़ी से बढ़ता है
जितना $L\omega$, और हर फेरा प्रतिरोध जोड़ देता है; लिट्ज़ तार $\delta$ से पतली बहुत
सारी विद्युतरोधी लड़ियाँ बाँध देता है, जिससे पूरा ताँबा चालन करता है।
\end{solution}

\begin{solution}{exo:b2:waves-in-media:6}
$\delta_{\text{Al}} = 11.9$, $0.84$, \qty{0.084}{mm}: $d/\delta = 0.08$, $1.2$, $12$: $0.7$,
$10$, \qty{103}{dB}। \qty{50}{Hz} पर चालक पारदर्शी है; ऊँचे
$\mu_r$ का लोहा फ्लक्स मोड़ लेता है (चुंबकीय परिरक्षण) और $\delta$ को
$\sqrt{\mu_r}$ से घटा भी देता है।
\end{solution}

\begin{solution}{exo:b2:waves-in-media:7}
(क) $v_0 = eE_0/m\omega$; $\tfrac12nm\langle v^2\rangle = ne^2E_0^2/4m\omega^2 = \tfrac14\varepsilon_0
E_0^2\,\omega_p^2/\omega^2$। (ख) $\tfrac14\varepsilon_0E_0^2$ और $\tfrac14\varepsilon_0E_0^2(1 - \omega_p^2/
\omega^2)$। (ग) योग $\tfrac12\varepsilon_0E_0^2$ है (कोष्ठक $1$ है); $\langle\Pi
\rangle = E_0^2k/2\mu_0\omega$, और अनुपात $c^2k/\omega = v_g$ है। (घ) $\tfrac14\cdot\tfrac14
/\tfrac12 = 1/8$।
\end{solution}

\begin{solution}{exo:b2:waves-in-media:8}
(क) $1/v_g = (1/c)(1 - \omega_p^2/\omega^2)^{-1/2} \approx (1 + \omega_p^2/2\omega^2)/c$। (ख) $\Delta t
= (L\omega_p^2/2c)(1/\omega_1^2 - 1/\omega_2^2)$ जहाँ $\omega_p^2 = ne^2/\varepsilon_0m$ और $\omega =
2\pi f$: अचर $e^2/8\pi^2\varepsilon_0mc = \qty{1.34e-7}{}$ (SI)। (ग) $nL = \qty{9.3e23}{m^{-2}}$,
$(1/f_1^2 - 1/f_2^2) = \qty{5.7e-18}{s^2}$: $\Delta t = \qty{0.72}{s}$। (घ) स्तंभ $nL$
(परिक्षेपण माप), और इससे $n$ के किसी प्रतिरूप के लिए दूरी।
\end{solution}

\begin{solution}{exo:b2:waves-in-media:9}
(क) $\operatorname{Im}\underline\chi = \omega_p^2\Gamma\omega/[(\omega_0^2 - \omega^2)^2 + \Gamma^2\omega^2]$, जिसमें
$\omega_0^2 - \omega^2 \approx 2\omega_0(\omega_0 - \omega)$: $n'' = \tfrac12\operatorname{Im}\underline\chi$ वही
बताया गया लोरेंत्सी है। (ख) $n''(\omega_0) = \omega_p^2/2\omega_0\Gamma$, $\alpha = 2n''\omega_0/c =
\omega_p^2/\Gamma c$। (ग) $\omega_p^2 = \qty{3.2e20}{s^{-2}}$: $\alpha = \qty{1.8e4}{m^{-1}}$;
$\ln10/\alpha = \qty{0.13}{mm}$। (घ) $n' - 1 \propto (\omega_0 - \omega)/[(\omega_0 - \omega)^2 +
\Gamma^2/4]$: नीचे धनात्मक, ऊपर ऋणात्मक, और रेखा के $\pm\Gamma/2$ के
भीतर घटता हुआ --- वहीं विषम।
\end{solution}

\begin{solution}{exo:b2:waves-in-media:10}
(क) $\underline E = E_0\eu^{-x/\delta}\eu^{\iu(\omega t - x/\delta)}$, $\underline B = \underline k\underline E/\omega
= (1 - \iu)\underline E/\omega\delta$। (ख) $\langle\Pi_x\rangle = \tfrac12\operatorname{Re}(\underline E
\underline B^*)/\mu_0 = E_0^2/2\mu_0\omega\delta$, $x = 0$ पर। (ग) $\int_0^\infty\tfrac12\gamma E_0^2
\eu^{-2x/\delta}\dd x = \gamma E_0^2\delta/4 = E_0^2/2\mu_0\omega\delta$, $\gamma = 2/\mu_0\omega\delta^2$ लेकर।
(घ) $H_0 = \sqrt2E_0/\mu_0\omega\delta$, अतः $\tfrac12R_sH_0^2 = E_0^2/\gamma\mu_0^2\omega^2\delta^3 =
E_0^2/2\mu_0\omega\delta$। \qty{1}{GHz} पर ताँबा: $R_s = 1/\gamma\delta = \qty{8}{m\ohm}$;
$\tfrac12 \times 8 \times 10^{-3} \times 10^4 = \qty{40}{W/m^2}$।
\end{solution}

\begin{solution}{exo:b2:waves-in-media:11}
(क) $\omega_p = \qty{1.4e16}{rad/s}$, $\lambda_p = 2\pi c/\omega_p = \qty{140}{nm}$। (ख) $\omega = \qty{3.8e15}{rad/s}
< \omega_p$: $1/\kappa = c/\sqrt{\omega_p^2 - \omega^2} = \qty{23}{nm}$, पूर्ण परावर्तन। (ग)
अब भी $\omega_p$ से नीचे: परावर्ती; पारदर्शिता केवल \qty{140}{nm} के परे (सोडियम में
\qty{210}{nm})। (घ) टक्करें $\underline\gamma$ को कोई वास्तविक भाग और कुछ प्रतिशत
अवशोषण दे देती हैं; और दृश्य में पड़ते बद्ध-इलेक्ट्रॉन अनुनाद
(सोना, ताँबा) नीले को सोखकर धातु को रंग दे देते हैं।
\end{solution}

\begin{solution}{exo:b2:waves-in-media:12}
(क) $\omega\tau = 0.12$: $\underline\varepsilon_r = 84 - 9.7\iu$; \qty{20}{GHz} पर $\omega\tau = 1$:
$45 - 40\iu$। (ख) $\underline n \approx 9.2 - 0.53\iu$; $c/2n''\omega = \qty{1.8}{cm}$: भट्ठी
बाहरी दो सेंटीमीटर पकाती है, बाक़ी चालन कर देता है। (ग) $\tfrac12 \times
1.54 \times 10^{10} \times 8.85 \times 10^{-12} \times 9.7 \times 4 \times 10^6 = \qty{2.6}{MW/m^3}$। (घ)
\qty{20}{GHz} पर अवशोषण लंबाई कोई मिलीमीटर होती: केवल त्वचा
पकती; \qty{2.45}{GHz} भीतर घुसता है, और वह कोई मुक्त पट्टी है।
\end{solution}

\begin{solution}{pb:b2:waves-in-media:1}
\textbf{1.} $f_p = \qty{9}{MHz}$: GPS पार जाता है, लघुतरंग परावर्तित हो जाती है।

\textbf{2.} $k = (\omega/c)\sqrt{1 - \omega_p^2/\omega^2}$: $v_\varphi = \omega/k \approx c(1 + \omega_p^2/
2\omega^2)$, $v_g = \dd\omega/\dd k = c^2k/\omega \approx c(1 - \omega_p^2/2\omega^2)$।

\textbf{3.} $\Delta t = \int(1/v_g - 1/c)\dd z = \int\omega_p^2\dd z/2c\omega^2 = e^2N_T/2\varepsilon_0
mc\omega^2 = (e^2/8\pi^2\varepsilon_0mc)N_T/f^2 = 1.34 \times 10^{-7} \times 2 \times 10^{17}/2.48 \times
10^{18} = \qty{11}{ns}$: अर्थात् \qty{3.2}{m}।

\textbf{4.} $N_T = (\Delta t_1 - \Delta t_2)/K(1/f_1^2 - 1/f_2^2)$; अंतर
\qty{7}{ns} है: $1\%$ के लिए \qty{0.07}{ns} चाहिए।

\textbf{5.} $v_\varphi$ $c$ से उतना ही ऊपर जाता है जितना $v_g$ नीचे रह जाता है:
वाहक की कला जल्दी पहुँचती है जबकि मॉडुलन देर से। कोई
वाहक-कला अभिग्राही संशोधन उलटे चिह्न से लगाता है।

\textbf{6.} दिन में \qty{3.2}{m}, रात में \qty{0.3}{m}।

\textbf{7.} टक्करों के साथ $f_p = \qty{0.9}{MHz}$: मध्यम और लघु तरंगें
सोख ली जाती हैं (कोई रेडियो अंधकार); \qty{1.5}{GHz} पर GPS को
$\propto 1/\omega^2$ अवशोषण मिलता है, जो नगण्य है।

\textbf{8.} $f_1\Delta t = 1.575 \times 10^9 \times 1.1 \times 10^{-8} = 17$ चक्र।

\textbf{9.} $\delta = \sqrt{2/\mu_0\gamma\omega} = \qty{10}{\micro m}$: कोई मिलीमीटर की दीवार
सौ त्वचा गहराइयाँ है।

\textbf{10.} $R_s = 1/\gamma\delta = \qty{0.1}{\ohm}$; $\tfrac12R_sH_0^2 = \qty{45}{W/m^2}$: \qty{45}{W},
$6\%$।

\textbf{11.} $\eu^{-\pi \times 0.5} = 0.21$: आयाम में \qty{-14}{dB}, शक्ति में
\qty{-27}{dB} (और नन्हा छेद वैसे भी बहुत कम युग्मित होता है)।

\textbf{12.} प्रकाश, $\lambda = \qty{0.5}{\micro m}$, छेदों से दो हज़ार गुना छोटा
है और पार निकल जाता है; सूक्ष्मतरंग, $\lambda = \qty{12}{cm}$, सौ गुना
बड़ी है।

\textbf{13.} $\underline n = 2 - 0.005\iu$; $c/2n''\omega = \qty{1.9}{m}$: काँच मुश्किल से
गरम होता है।

\textbf{14.} किसी दस-माइक्रोमीटर परत में $j \approx H_0/\delta = \qty{3e7}{A/m^2}$,
$j^2/\gamma = \qty{9e8}{W/m^3}$: नोकें मिलीसेकंडों में तपती हैं, इलेक्ट्रॉन उगलती हैं,
हवा को आयनित कर देती हैं --- चिंगारियाँ।

\textbf{15.} $\delta = \qty{18}{\micro m}$; अनुप्रस्थ काट $2\pi a\delta = \qty{5.6e-8}{m^2}$ के सामने
$\pi a^2 = \qty{7.9e-7}{m^2}$: $R_{\text{प्रत्यावर्ती}} = \qty{0.30}{\ohm/m}$, अर्थात् दिष्ट का चौदह गुना।

\textbf{16.} $RI^2 = \qty{120}{W/m}$; प्रति मीटर सतह $\pi d = \qty{3.1e-3}{m^2}$:
$\Delta T \approx \qty{1900}{K}$ --- वह पिघल जाता; ऐसी कुंडलियाँ जल-शीतित नल होती हैं।

\textbf{17.} $\delta = \qty{0.14}{mm}$: प्रेरित धारा, और गर्मी, दसवें
मिलीमीटर में ही रह जाती है; कोई छोटी स्पंद और उसके बाद कोई शमन गियर की
सतह कठोर कर देता है जबकि भीतर का भाग सख़्तजान बना रहता है।

\textbf{18.} $R = 1/\gamma2\pi a\delta = \qty{0.10}{\ohm/m}$, अर्थात् तीन गुना कम; भीतरी
भाग कुछ नहीं ले चलता --- इसीलिए कोई नल, जहाँ ताँबा बेकार जाता वहाँ
शीतक भर दिया गया।

\textbf{19.} $\delta = a$ तब जब $f = 1/\pi\mu_0\gamma a^2 = \qty{17}{Hz}$; उससे नीचे धारा
एकसमान है और दिष्ट सूत्र चलता है।

\textbf{20.} $n^2 \approx 1 + \omega_p^2/\omega_0^2 + (\omega_p^2/\omega_0^2)(\lambda_0/\lambda)^2$: कौशी का रूप,
जिसमें $n(\infty)^2 - 1 = \omega_p^2/\omega_0^2 = 1.25$ जब $n(\infty) = 1.5$।

\textbf{21.} $n = 1.531$ और $1.510$, $v = 1.96$ और \qty{1.99e8}{m/s}; $n_g(550) =
1.550$; $n_g(400) - n_g(700) = 0.063$: किलोमीटर भर पर \qty{210}{ns}।

\textbf{22.} $\sin((A + D)/2) = 1.5165 \times 0.5$: $D = \qty{38.6}{\degree}$; $\dd D = 2\sin
(A/2)\dd n/\cos((A + D)/2) = 1.53\,\dd n$: $\Delta D = \qty{1.85}{\degree}$, \qty{2}{m} पर \qty{6.5}{cm}।

\textbf{23.} \qty{200}{nm} पर आवृत्ति अनुनाद के पास पहुँच जाती है और
$n''$ बढ़ जाता है: अपारदर्शी। \qty{700}{nm} पर कोई दुर्बल अनुनाद लाल सोख लेता है:
हरा काँच।

\textbf{24.} नीला पराबैंगनी अनुनाद के अधिक पास है: बद्ध
इलेक्ट्रॉन अधिक उत्तर देते हैं, $n$ बड़ा होता है, और विचलन भी।

\textbf{25.} \qty{1.5}{GHz} पर F परत: काल्पनिक $\underline\gamma$, और किसी छोटे
समूह विलंब के साथ संचरण। \qty{2.45}{GHz} पर इस्पात: वास्तविक $\underline\gamma$, \omterm{thm:b2:waves-in-media:skin}{त्वचा
प्रभाव}, परावर्तन। \qty{13.56}{MHz} पर ताँबा: वही, और \qty{18}{\micro m} की
कोई त्वचा। दृश्य में काँच: $\iu\omega\varepsilon_0\chi$, कोई वास्तविक अपवर्तनांक, कोई
पारदर्शी \omterm{def:b2:dispersion-wave-packets:relation}{परिक्षेपी माध्यम}।
\end{solution}
